% Authoritative LaTeX chapter. Edit this file for future revisions.
% Initially migrated 2026-09-11; no Markdown import occurs during PDF builds.
\chapter{Basic Concepts and Observation Geometry}
\label{chap:01_concepts_geometry}
Manual contents · \hyperref[chap:05_output_formats]{Output formats}

This document is based on the public inputs and computed outputs of \texttt{code/\allowbreak{}OCRT\_\allowbreak{}C/\allowbreak{}src}. The reference is Git commit \texttt{54861c68e35ec5aa\allowbreak{}a3938fad53e93dc2\allowbreak{}dfa7eb1d}; the executable's version string alone does not establish whether it reflects the latest source. Historical validation reports document the versions used at the time.

\section{What OCRT calculates}
\label{sec:auto:01_concepts_geometry:1}
OCRT is a radiative transfer model that calculates the scattering, absorption, reflection, and transmission of sunlight in the atmosphere, at the sea surface, and underwater. It computes Stokes \texttt{I}, \texttt{Q}, and \texttt{U} for elastic scattering at the same wavelength. \texttt{I} represents total radiance, while \texttt{Q} and \texttt{U} are the linear polarization components relative to the chosen reference plane. Circular polarization \texttt{V} is not included in the public output.

Inputs describe observation geometry, wavelength, atmospheric pressure, aerosols and gases, wind at the sea surface, and the optical properties of water. The main outputs are as follows.


{\TableFont
\begin{longtable}{@{}>{\raggedright\arraybackslash}p{\dimexpr 0.3400\linewidth-4.00pt\relax}>{\raggedright\arraybackslash}p{\dimexpr 0.6600\linewidth-4.00pt\relax}@{}}
\toprule
\textbf{Location or quantity} & \textbf{Meaning} \\
\midrule
\endfirsthead
\toprule
\textbf{Location or quantity} & \textbf{Meaning} \\
\midrule
\endhead
\bottomrule
\endfoot
TOA & Signal leaving the top of the atmosphere toward the satellite \\
\texttt{0+} & Immediately above the water surface, on the air side \\
\texttt{0-\allowbreak{}} & Immediately below the water surface, on the water side \\
\texttt{rho} & Directional TOA reflectance normalized by incident solar illumination; dimensionless \\
\texttt{Rrs} & \texttt{Lu(0+)/\allowbreak{}Ed(0+)}, above-surface remote-sensing reflectance, sr\textsuperscript{-}\textsuperscript{1} \\
\texttt{rrs} & \texttt{Lu(0-\allowbreak{})/\allowbreak{}Ed(0-\allowbreak{})}, below-surface remote-sensing reflectance, sr\textsuperscript{-}\textsuperscript{1} \\
\texttt{Lu}, \texttt{Ed}, \texttt{Eu} & Upwelling radiance, downwelling irradiance, and upwelling irradiance, respectively \\
\end{longtable}
}

\texttt{Rrs} and \texttt{rho} have different denominators and units. Surface reflectance obtained by multiplying \texttt{Rrs} by \ensuremath{\pi} is also different from TOA \texttt{rho}, which includes the atmospheric path. Raw \texttt{Lu} and \texttt{Ed} values use model-specific incident-light normalization, so first consult the \hyperref[chap:05_output_formats]{output normalization rules}.

The calculation assumes horizontally uniform layers. Multiple scattering is calculated by SOS (Successive Orders of Scattering, accumulated by scattering order), and azimuthal dependence is represented by Fourier modes. Enabling \texttt{-\allowbreak{}-\allowbreak{}pssa} applies the IPSS spherical correction in the current implementation. This is not a model for fully three-dimensional terrain/cloud radiative transfer or for arbitrary seabed topography.

\section{Select the surface mode first}
\label{sec:auto:01_concepts_geometry:2}

{\TableFont
\begin{longtable}{@{}>{\raggedright\arraybackslash}p{\dimexpr 0.3000\linewidth-5.33pt\relax}>{\raggedright\arraybackslash}p{\dimexpr 0.2300\linewidth-5.33pt\relax}>{\raggedright\arraybackslash}p{\dimexpr 0.4700\linewidth-5.33pt\relax}@{}}
\toprule
\textbf{\texttt{-\allowbreak{}-\allowbreak{}surface}} & \textbf{Calculation domain} & \textbf{Typical use} \\
\midrule
\endfirsthead
\toprule
\textbf{\texttt{-\allowbreak{}-\allowbreak{}surface}} & \textbf{Calculation domain} & \textbf{Typical use} \\
\midrule
\endhead
\bottomrule
\endfoot
\texttt{black} & Atmosphere above a nonreflecting lower boundary & Pure atmospheric path; checking Rayleigh scattering \\
\texttt{flat} & Atmosphere above a flat Fresnel-reflecting boundary & Reference calculations for flat-surface Fresnel reflection \\
\texttt{black\_\allowbreak{}fresnel\_\allowbreak{}ocean} & Atmosphere and a wind-dependent Fresnel sea surface, with no water-leaving light & Rayleigh/aerosol path LUTs for ocean atmospheric correction \\
\texttt{ocean} & Coupled atmosphere--sea-surface--underwater radiative transfer & \texttt{Rrs}, \texttt{rrs}, TOA ocean signals and simulated data \\
\end{longtable}
}

For \texttt{ocean}, specify water properties with one of \texttt{-\allowbreak{}-\allowbreak{}water-\allowbreak{}model ocrt}, \texttt{ccrr}, or \texttt{iop}. \texttt{ocrt} and \texttt{ccrr} select how constituents are converted into optical properties; \texttt{iop} supplies absorption, scattering, and backscattering coefficients and the phase function directly. Even identically named concentrations can produce different optical properties under different constituent models.

For IOPs (Inherent Optical Properties), \texttt{a}, \texttt{b}, and \texttt{bb} are the absorption, scattering, and backscattering coefficients, respectively, in m\textsuperscript{-}\textsuperscript{1}. The extinction coefficient is \texttt{c=\allowbreak{}a+b}, and the single-scattering albedo is \texttt{omega=\allowbreak{}b/\allowbreak{}(a+b)}. Since \texttt{bb} is the backward-hemisphere component of \texttt{b}, do not use \texttt{a+b+bb} as the extinction coefficient.

\section{SZA, VZA, and RAA: always referenced to the ground pixel}
\label{sec:auto:01_concepts_geometry:3}

{\TableFont
\begin{longtable}{@{}>{\raggedright\arraybackslash}p{\dimexpr 0.3000\linewidth-5.33pt\relax}>{\raggedright\arraybackslash}p{\dimexpr 0.2300\linewidth-5.33pt\relax}>{\raggedright\arraybackslash}p{\dimexpr 0.4700\linewidth-5.33pt\relax}@{}}
\toprule
\textbf{Input / CSV name} & \textbf{Definition} & \textbf{Unit} \\
\midrule
\endfirsthead
\toprule
\textbf{Input / CSV name} & \textbf{Definition} & \textbf{Unit} \\
\midrule
\endhead
\bottomrule
\endfoot
\texttt{-\allowbreak{}-\allowbreak{}sza}, \texttt{sza\_\allowbreak{}deg} & Angle between the upward vertical at the ground pixel and the direction toward the Sun & Degrees \\
\texttt{-\allowbreak{}-\allowbreak{}vza}, \texttt{vza\_\allowbreak{}deg} & Angle between the upward vertical at the ground pixel and the direction toward the sensor & Degrees \\
\texttt{-\allowbreak{}-\allowbreak{}raa}, \texttt{raa\_\allowbreak{}deg} & OCRT public relative azimuth angle, following the operational definition below & Degrees \\
\end{longtable}
}

\texttt{VZA=\allowbreak{}0\textdegree{}} is nadir viewing. Distinguish VZA from the off-nadir angle viewed downward from the satellite: in spherical geometry, these angles can differ. Inputs remain referenced to the ground pixel when \texttt{-\allowbreak{}-\allowbreak{}pssa} is used. For \texttt{ocean}, VZA is the viewing angle in air, and the standard path obtains the underwater angle through an internal refraction transformation. Do not reinterpret the same numerical VZA as an underwater angle.

For daytime upward viewing, generally prepare valid geometries with \texttt{0 \ensuremath{\leq} SZA < 90} and \texttt{0 \ensuremath{\leq} VZA < 90}. The default LUT spans VZA \texttt{0--85\textdegree{}}; inspect the coordinates in the output CSV for the exact grid. Normalize input RAA to \texttt{[0,\allowbreak{}360)} where possible. \texttt{360\textdegree{}} is physically the same direction as \texttt{0\textdegree{}} and is not included in the LUT as a duplicate endpoint.

Figure~\ref{fig:geometry-zenith-refraction} distinguishes the two zenith angles from underwater refraction. Figure~\ref{fig:geometry-curvature} shows why satellite off-nadir angle must not be passed directly to \texttt{--vza}.

% Editable vector artwork. Both editions use the same physical geometry.
\begin{figure}[H]
\centering
\begin{tikzpicture}[x=1cm,y=1cm,font=\sffamily\fontsize{9}{11}\selectfont,
  ray/.style={line width=1.05pt,-{Latex[length=2.2mm]}},
  axis/.style={ocrtGray,dashed,line width=.55pt},
  angle/.style={ocrtNavy,line width=.8pt},
  note/.style={align=center,text width=6.8cm}]
\path[use as bounding box] (0,-3.6) rectangle (15.8,3.9);
\begin{scope}[shift={(3.6,0)}]
  \node[anchor=north,font=\sffamily\bfseries] at (0,3.9)
    {\FigLang{(a) 화소에서 위쪽을 바라본 각도}{(a) Look directions from the pixel}};
  \fill[ocrtTeal!8] (-3.3,-.55) rectangle (3.3,0);
  \draw[ocrtTeal,line width=.8pt] (-3.3,0)--(3.3,0);
  \draw[axis] (0,-.35)--(0,2.75);
  \node[anchor=south,fill=white,inner sep=1.5pt] at (0,2.75)
    {\FigLang{국소 수직선}{Local vertical}};
  \draw[ocrtGray,dashed,-{Latex[length=2mm]}] (0,0)--(125:2.65);
  \draw[ray,orange!85!black] (125:2.25)--(125:.35);
  \draw[ray,ocrtTeal] (0,0)--(45:2.8);
  \node[anchor=south,align=center] at (-1.85,2.37)
    {\FigLang{태양을 향한 방향}{Toward the Sun}};
  \node[anchor=south,align=center] at (2.02,2.08)
    {\FigLang{센서를 향한 방향}{Toward the sensor}};
  \draw[angle] (90:.85) arc[start angle=90,end angle=125,radius=.85];
  \node[ocrtNavy] at (-.52,1.14) {SZA};
  \draw[angle] (45:1.12) arc[start angle=45,end angle=90,radius=1.12];
  \node[ocrtNavy] at (.66,1.37) {VZA};
  \fill[ocrtNavy] (0,0) circle (2pt);
  \node[anchor=north] at (0,-.12) {\FigLang{지상 화소}{Ground pixel}};
  \node[note,anchor=north] at (0,-.95)
    {\FigLang{태양을 바라보는 방향과 입사광의 진행 방향은 서로 반대다.}{The Sunward look direction is opposite to incoming photon travel.}};
  \draw[ray,orange!85!black] (-2.75,-2.13)--(-1.9,-2.13);
  \node[anchor=west] at (-1.7,-2.13) {\FigLang{입사 태양광}{Incoming sunlight}};
  \draw[ray,ocrtTeal] (-2.75,-2.68)--(-1.9,-2.68);
  \node[anchor=west] at (-1.7,-2.68) {\FigLang{센서로 나가는 빛}{Light toward the sensor}};
\end{scope}
\draw[ocrtLine] (7.8,-3.35)--(7.8,3.55);
\begin{scope}[shift={(11.8,0)}]
  \node[anchor=north,font=\sffamily\bfseries] at (0,3.9)
    {\FigLang{(b) 공기 중 VZA와 수중 각도}{(b) Air VZA and underwater angle}};
  \fill[ocrtTeal!8] (-3.4,-1.9) rectangle (3.4,0);
  \draw[ocrtTeal,line width=.8pt] (-3.4,0)--(3.4,0);
  \draw[axis] (0,-1.8)--(0,2.65);
  \draw[ray,ocrtTeal] (-.847,-1.55)--(0,0);
  \draw[ray,ocrtTeal] (0,0)--(50:2.65);
  \node[anchor=south,align=center] at (1.9,2.1)
    {\FigLang{센서}{Sensor}};
  \node[anchor=west] at (-3.18,.42) {\FigLang{공기}{Air}};
  \node[anchor=west] at (-3.18,-.43) {\FigLang{물}{Water}};
  \draw[angle] (50:.9) arc[start angle=50,end angle=90,radius=.9];
  \node[ocrtNavy] at (.7,1.17) {$\mathrm{VZA}_{\rm air}$};
  \draw[angle] (-90:.7) arc[start angle=-90,end angle=-118.66,radius=.7];
  \node[ocrtNavy] at (-.63,-1.05) {$\theta_w$};
  \fill[ocrtNavy] (0,0) circle (2pt);
  \node[note,anchor=north] at (0,-2.16)
    {$n_{\rm air}\sin(\mathrm{VZA}_{\rm air})=n_w\sin\theta_w$};
  \node[note,anchor=north] at (0,-2.7)
    {\FigLang{입력: 공기 중 VZA\\수중 각도는 내부에서 계산}{Input: air VZA\\Underwater angle is computed internally}};
\end{scope}
\end{tikzpicture}
\caption{\FigLang{SZA와 VZA는 화소의 위쪽 수직선에서 잰다. (a)는 두 look 방향을 한 단면에 배치한 예시이며 일반 기하의 RAA는 별도로 지정한다. (b)는 평탄한 경계의 굴절 도식이다. 표준 해양 경로의 VZA 입력을 수중 각도로 바꾸어 넣지 않는다.}{SZA and VZA are measured from the upward vertical at the pixel. Panel (a) places both look directions in one cross-section for illustration; RAA is specified separately for general geometry. Panel (b) shows refraction at a flat interface. Do not replace the standard ocean-path air-VZA input with an underwater angle.}}
\label{fig:geometry-zenith-refraction}
\end{figure}

% Angles drawn from exact ray endpoints; exaggerated orbital height is schematic.
\begin{figure}[H]
\centering
\begin{tikzpicture}[x=1cm,y=1cm,font=\sffamily\fontsize{9}{11}\selectfont,
  note/.style={align=left,text width=5.4cm}]
\path[use as bounding box] (0,-3.15) rectangle (15.8,4.72);
\coordinate (C) at (3.6,-2.5);
\coordinate (P) at (3.6,1);
\coordinate (S) at (8.6,4);
\fill[ocrtTeal!9] (C)--($(C)+(25:3.5)$)
  arc[start angle=25,end angle=155,radius=3.5]--cycle;
\draw[ocrtTeal,line width=1pt] ($(C)+(25:3.5)$)
  arc[start angle=25,end angle=155,radius=3.5];
\draw[ocrtLine,dotted,line width=.8pt] ($(C)+(25:3.85)$)
  arc[start angle=25,end angle=155,radius=3.85];
\draw[ocrtGray,dashed] (C)--(3.6,4.25);
\draw[ocrtGray,dashed] (S)--(C);
\draw[ocrtTeal,line width=1.15pt,-{Latex[length=2.2mm]}] (P)--(S);
\draw[ocrtNavy,line width=.8pt] ($(P)+(30.96376:1)$)
  arc[start angle=30.96376,end angle=90,radius=1];
\node[ocrtNavy,fill=white,inner sep=1pt] at (4.24,2.36) {VZA};
\draw[purple!70!black,line width=.8pt] ($(S)+(-149.03624:1.1)$)
  arc[start angle=-149.03624,end angle=-127.56859,radius=1.1];
\node[purple!70!black,fill=white,inner sep=1pt] at (7.25,2.87) {$\beta$};
\fill[ocrtNavy] (P) circle (2pt);
\fill[ocrtNavy] (C) circle (2pt);
\fill[ocrtTeal] (S) circle (2.3pt);
\node[anchor=south,align=center] at (3.6,4.29)
  {\FigLang{화소의 국소 수직선}{Local vertical at the pixel}};
\node[anchor=south,align=center] at (8.45,4.19)
  {\FigLang{위성}{Satellite}};
\node[align=center] at (2.3,1.95)
  {\FigLang{지상 화소}{Ground pixel}};
\draw[ocrtGray,line width=.5pt] (2.8,1.66)--(3.5,1.1);
\node[anchor=north,align=center] at (3.6,-2.69)
  {\FigLang{지구 중심}{Earth center}};
\node[align=center] at (1.3,.95)
  {\FigLang{대기층}{Atmosphere}};
\draw[ocrtGray,line width=.5pt] (2.22,1.02)--(2.5,1.189);
\node[fill=ocrtTeal!9,inner sep=2pt] at (3.22,-.65) {$R$};
\node[rotate=52.43,fill=white,inner sep=2pt] at (5.8,.47) {$R+h$};
\node[anchor=north west,note,font=\sffamily\bfseries] at (10.0,3.9)
  {\FigLang{입력하는 각도는 VZA}{Enter VZA at the pixel}};
\node[anchor=north west,note] at (10.0,3.1)
  {\FigLang{\texttt{--vza}: 화소의 위쪽 수직선과 센서 look 방향 사이의 각도.}{\texttt{--vza}: angle between the pixel's upward vertical and its sensor look direction.}};
\node[anchor=north west,note] at (10.0,1.65)
  {\FigLang{$\beta$: 위성에서 지구 중심을 향하는 천저 방향과 화소를 향하는 방향 사이의 off-nadir 각도.}{$\beta$: satellite off-nadir angle, between the Earth-center direction and the look toward the pixel.}};
\node[anchor=north west,note] at (10.0,-.05)
  {$R\sin(\mathrm{VZA})=(R+h)\sin\beta$};
\node[anchor=north west,note] at (10.0,-.76)
  {\FigLang{$R$: 구의 반지름, $h$: 센서 고도. 직선 광선·구형 지구에서의 기하 관계다.}{$R$: sphere radius; $h$: sensor altitude. This relation assumes straight rays and a spherical Earth.}};
\node[anchor=north west,note] at (10.0,-2.24)
  {\FigLang{\texttt{--pssa}에서도 입력 기준은 화소다.}{The input remains pixel-referenced with \texttt{--pssa}.}};
\end{tikzpicture}
\caption{\FigLang{구면 기하에서 지상 화소의 VZA와 위성의 off-nadir 각도는 일반적으로 다르다. 대기층 두께와 센서 고도는 구별이 쉽도록 과장했으며 비례도·계산 결과가 아니다. 현재 \texttt{--pssa}의 IPSS 보정도 지상 화소 기준 기하를 사용한다. 이 도식은 완전한 3차원 구면 다중산란 모델을 의미하지 않는다.}{In spherical geometry, ground-pixel VZA and satellite off-nadir angle generally differ. Atmospheric thickness and sensor altitude are exaggerated for clarity; this is not a scale drawing or a computed result. The current IPSS correction enabled by \texttt{--pssa} also uses ground-pixel geometry. The schematic does not imply a fully three-dimensional spherical multiple-scattering model.}}
\label{fig:geometry-curvature}
\end{figure}


\section{Operational definition of OCRT RAA}
\label{sec:auto:01_concepts_geometry:4}
\textbf{The direct sunglint direction in OCRT is \texttt{RAA=\allowbreak{}180\textdegree{}}.} For specular reflection from a flat surface, this is the specular branch when \texttt{SZA=\allowbreak{}VZA}. Wind broadens the specular signal into neighboring angles through the Cox--Munk slope distribution.


{\TableFont
\begin{longtable}{@{}>{\raggedright\arraybackslash}p{\dimexpr 0.3400\linewidth-4.00pt\relax}>{\raggedright\arraybackslash}p{\dimexpr 0.6600\linewidth-4.00pt\relax}@{}}
\toprule
\textbf{OCRT RAA} & \textbf{Interpretation} \\
\midrule
\endfirsthead
\toprule
\textbf{OCRT RAA} & \textbf{Interpretation} \\
\midrule
\endhead
\bottomrule
\endfoot
\texttt{0\textdegree{}} & Principal-plane branch opposite direct glint \\
\texttt{90\textdegree{}} & One azimuth perpendicular to the principal plane \\
\texttt{180\textdegree{}} & Direct-glint principal-plane branch; specular geometry when \texttt{SZA=\allowbreak{}VZA} \\
\texttt{270\textdegree{}} & Mirror azimuth of \texttt{90\textdegree{}}; \texttt{I/\allowbreak{}Q} are symmetric, while \texttt{U} can have the opposite sign \\
\texttt{360\textdegree{}} & Same direction as \texttt{0\textdegree{}} \\
\end{longtable}
}

Do not determine RAA from descriptions such as “on the same side as the Sun,” “forward,” or “backward” alone. The direction \textbf{looking toward the Sun} and the direction in which sunlight \textbf{propagates} differ by 180\textdegree{} in horizontal azimuth. If another radiative transfer code or LUT places glint at \texttt{RAA=\allowbreak{}0\textdegree{}}, that numerical value cannot be used directly in OCRT.

Following solar photon travel and the sensor look direction in Figure~\ref{fig:geometry-glint-branches} makes the two branches easier to distinguish. The \texttt{RAA=180\textdegree{}} branch alone does not define the specular center: \texttt{SZA=VZA} is also required.

% Equal 30-degree zenith angles. Solar-look and propagation vectors differ.
\begin{figure}[H]
\centering
\begin{tikzpicture}[x=1cm,y=1cm,font=\sffamily\fontsize{9}{11}\selectfont,
  ray/.style={line width=1.15pt,-{Latex[length=2.2mm]}},
  note/.style={align=center,text width=6.9cm}]
\path[use as bounding box] (0,-3.15) rectangle (15.8,4.25);
\begin{scope}[shift={(3.6,0)}]
  \node[anchor=north,font=\sffamily\bfseries] at (0,4.25)
    {(a) RAA $=0^\circ$};
  \fill[ocrtTeal!8] (-3.3,-.4) rectangle (3.3,0);
  \draw[ocrtTeal,line width=.8pt] (-3.3,0)--(3.3,0);
  \draw[ocrtGray,dashed] (0,0)--(0,2.45);
  \draw[ocrtGray,dashed] (0,0)--(120:3.0);
  \draw[ray,orange!85!black] (120:2.65)--(120:1.7);
  \draw[ray,ocrtTeal] (120:.15)--(120:1.35);
  \node[anchor=south,align=center,text width=4.8cm] at (-.9,2.75)
    {\FigLang{태양과 센서의 look 방향 일치}{Solar and sensor look directions coincide}};
  \draw[ocrtNavy] (90:.8) arc[start angle=90,end angle=120,radius=.8];
  \node[anchor=west,align=left] at (.3,1.25)
    {SZA $=$ VZA $=30^\circ$};
  \fill[ocrtNavy] (0,0) circle (2pt);
  \node[anchor=north] at (0,-.12) {\FigLang{화소}{Pixel}};
  \node[note,anchor=north] at (0,-.79)
    {\FigLang{평탄면의 직접 글린트 가지가 아님}{Outside the flat-surface direct-glint branch}};
  \node[note,anchor=north] at (0,-1.88)
    {$\Theta=180^\circ$\\\FigLang{대기 입자의 정확 후방산란 기하}{Exact atmospheric backscattering geometry}};
\end{scope}
\draw[ocrtLine] (7.8,-2.95)--(7.8,3.85);
\begin{scope}[shift={(11.8,0)}]
  \node[anchor=north,font=\sffamily\bfseries] at (0,4.25)
    {(b) RAA $=180^\circ$};
  \fill[ocrtTeal!8] (-3.3,-.4) rectangle (3.3,0);
  \draw[ocrtTeal,line width=.8pt] (-3.3,0)--(3.3,0);
  \draw[ocrtGray,dashed] (0,0)--(0,2.85);
  \draw[ray,orange!85!black] (120:2.95)--(0,0);
  \draw[ray,ocrtTeal] (0,0)--(60:2.95);
  \node[anchor=south] at (-1.5,2.7) {\FigLang{태양}{Sun}};
  \node[anchor=south] at (1.5,2.7) {\FigLang{센서}{Sensor}};
  \draw[ocrtNavy] (90:.9) arc[start angle=90,end angle=120,radius=.9];
  \draw[ocrtNavy] (60:.9) arc[start angle=60,end angle=90,radius=.9];
  \node[ocrtNavy] at (-.55,1.16) {SZA};
  \node[ocrtNavy] at (.55,1.16) {VZA};
  \fill[ocrtNavy] (0,0) circle (2pt);
  \node[anchor=north] at (0,-.12) {\FigLang{화소}{Pixel}};
  \node[note,anchor=north] at (0,-.79)
    {SZA $=$ VZA $=30^\circ$\\\FigLang{평탄면 정반사}{Flat-surface specular reflection}};
  \node[note,anchor=north] at (0,-1.88)
    {$\Theta=120^\circ$\\\FigLang{수면 정반사와 입자 후방산란은 구별}{Surface specular reflection differs from particle backscatter}};
\end{scope}
\end{tikzpicture}
\caption{\FigLang{두 주평면 가지를 같은 천정각으로 비교한 단면도. 주황 화살표는 입사광, 청록 화살표는 센서 방향으로 나가는 빛이다. (a)의 두 화살표는 같은 직선 위에서 서로 반대 방향을 가리킨다. $\Theta$는 입사 태양광의 진행 방향과 상향 관측 방향 사이의 대기 산란각이며 수면 법선에서 잰 각도가 아니다. 바람이 있는 해수면의 글린트는 (b)의 중심 주변으로 넓어질 수 있다.}{Cross-sections comparing the two principal-plane branches at equal zenith angles. Orange arrows show incoming photons and teal arrows show light toward the sensor. The arrows in (a) point in opposite directions on the same line. $\Theta$ is the atmospheric scattering angle between incoming solar propagation and upward viewing, not an angle measured from the surface normal. A wind-roughened surface can broaden glint around the center shown in (b).}}
\label{fig:geometry-glint-branches}
\end{figure}


\subsection{Calculation from satellite SAA/VAA}
\label{sec:auto:01_concepts_geometry:4.1}
First check the following in the product metadata.

\begin{enumerate}
\item Verify that SAA is the azimuth \textbf{toward the Sun} from the ground pixel.
\item Verify that VAA is the azimuth \textbf{toward the sensor} from the ground pixel. If it is the azimuth from the sensor toward the surface, add \texttt{180\textdegree{}} to convert it to a ground-pixel reference.
\item Align the zero reference and positive rotation direction of SAA and VAA. Do not repeat this procedure on a value already supplied as a relative azimuth angle.
\end{enumerate}

When both values satisfy these look-azimuth definitions and use the same rotation direction, the following is an explicit conversion convention consistent with OCRT's \texttt{0/\allowbreak{}180\textdegree{}} branches.


\begin{ManualCode}
def ocrt_raa_from_look_azimuths(saa_deg, vaa_deg):
    return (vaa_deg - saa_deg) % 360.0
\end{ManualCode}


{\TableFont
\begin{longtable}{@{}>{\raggedright\arraybackslash}p{\dimexpr 0.2300\linewidth-6.00pt\relax}>{\raggedright\arraybackslash}p{\dimexpr 0.2400\linewidth-6.00pt\relax}>{\raggedright\arraybackslash}p{\dimexpr 0.2300\linewidth-6.00pt\relax}>{\raggedright\arraybackslash}p{\dimexpr 0.3000\linewidth-6.00pt\relax}@{}}
\toprule
\textbf{SAA} & \textbf{VAA} & \textbf{OCRT RAA from the equation above} & \textbf{Meaning} \\
\midrule
\endfirsthead
\toprule
\textbf{SAA} & \textbf{VAA} & \textbf{OCRT RAA from the equation above} & \textbf{Meaning} \\
\midrule
\endhead
\bottomrule
\endfoot
30\textdegree{} & 30\textdegree{} & 0\textdegree{} & Solar look azimuth and sensor look azimuth coincide \\
30\textdegree{} & 210\textdegree{} & 180\textdegree{} & Glint branch; the specular center also requires SZA=VZA \\
350\textdegree{} & 10\textdegree{} & 20\textdegree{} & Difference crossing the 0\textdegree{} boundary \\
10\textdegree{} & 350\textdegree{} & 340\textdegree{} & Mirror azimuth of the preceding row \\
\end{longtable}
}

This equation is a transformation derived in the coordinates declared above; it does not automatically align the Stokes reference plane of a particular satellite product. Reversing the product's positive azimuth rotation direction changes \texttt{RAA} to \texttt{360-RAA}, reversing the sign of \texttt{U} even when \texttt{I/\allowbreak{}Q} remain the same. For polarized data, record the original SAA/VAA, conversion equation, reference plane, and rotation direction together, and compare signed \texttt{U} at an off-principal-plane geometry. If the product's azimuth and Stokes definitions are unknown, agreement in \texttt{I/\allowbreak{}Q} alone does not validate the \texttt{U} transformation.

Figure~\ref{fig:geometry-azimuth} shows the rotation from SAA to VAA. The right panel rotates the same coordinates to the solar look azimuth and illustrates why mirror azimuths need to remain distinguishable in stored data.

% Clockwise is an explicitly declared illustrative coordinate convention.
\begin{figure}[H]
\centering
\begin{tikzpicture}[x=1cm,y=1cm,font=\sffamily\fontsize{9}{11}\selectfont,
  ray/.style={line width=1pt,-{Latex[length=2mm]}},
  arcarr/.style={line width=.85pt,-{Latex[length=1.7mm]}},
  note/.style={align=center,text width=6.9cm}]
\path[use as bounding box] (0,-3.1) rectangle (15.8,3.7);
\begin{scope}[shift={(3.6,0)}]
  \node[anchor=north,font=\sffamily\bfseries] at (0,3.7)
    {\FigLang{(a) 위에서 본 SAA, VAA, RAA}{(a) SAA, VAA, RAA viewed from above}};
  \draw[ocrtLine] (0,0) circle (2.05);
  \draw[ocrtGray,dashed,-{Latex[length=2mm]}] (0,0)--(0,2.65);
  \node[anchor=south,align=center] at (0,2.68)
    {\FigLang{공통 방위 영점}{Shared azimuth zero}};
  \draw[ray,orange!85!black] (0,0)--(60:2.45);
  \draw[ray,ocrtTeal] (0,0)--(-30:2.45);
  \node[anchor=west,align=left] at (1.4,2.15)
    {\FigLang{태양 look}{Solar look}\\SAA $=30^\circ$};
  \node[anchor=north,align=center] at (2.0,-1.47)
    {\FigLang{센서 look}{Sensor look}\\VAA $=120^\circ$};
  \draw[arcarr,orange!85!black] (90:.7) arc[start angle=90,end angle=60,radius=.7];
  \node[anchor=east,orange!85!black] at (-.16,.95) {SAA};
  \draw[arcarr,ocrtTeal] (90:1.03) arc[start angle=90,end angle=-30,radius=1.03];
  \node[ocrtTeal,fill=white,inner sep=1pt] at (.96,.75) {VAA};
  \draw[arcarr,ocrtNavy] (60:1.65) arc[start angle=60,end angle=-30,radius=1.65];
  \node[ocrtNavy,fill=white,inner sep=1.5pt] at (2.14,.24) {RAA $=90^\circ$};
  \fill[ocrtNavy] (0,0) circle (2pt);
  \node[anchor=north east] at (-.1,-.1) {\FigLang{화소}{Pixel}};
  \node[note,anchor=north] at (0,-2.28)
    {$\mathrm{RAA}=(\mathrm{VAA}-\mathrm{SAA})\bmod 360^\circ$};
  \node[note,anchor=north] at (0,-2.83)
    {\FigLang{이 그림의 양의 회전: 시계 방향}{Positive rotation in this figure: clockwise}};
\end{scope}
\draw[ocrtLine] (7.8,-2.95)--(7.8,3.35);
\begin{scope}[shift={(11.8,0)}]
  \node[anchor=north,font=\sffamily\bfseries] at (0,3.7)
    {\FigLang{(b) 태양 look을 위쪽으로 맞춘 좌표}{(b) Rotate the solar look to the top}};
  \draw[ocrtLine] (0,0) circle (1.72);
  \draw[ray,orange!85!black] (0,0)--(0,2.02);
  \draw[ray,ocrtTeal] (0,0)--(1.92,0);
  \draw[ray,ocrtTeal] (0,0)--(0,-1.87);
  \draw[ray,ocrtTeal] (0,0)--(-1.92,0);
  \draw[purple!70!black,dashed,line width=.8pt] (0,0)--(70:1.72);
  \draw[purple!70!black,dashed,line width=.8pt] (0,0)--(110:1.72);
  \node[purple!70!black,anchor=west] at (.7,1.69) {$20^\circ$};
  \node[purple!70!black,anchor=east] at (-.7,1.69) {$340^\circ$};
  \node[anchor=south,align=center] at (0,2.1)
    {$0^\circ\;(=360^\circ)$\\\FigLang{태양 look}{Solar look}};
  \node[anchor=west] at (2.05,0) {$90^\circ$};
  \node[anchor=east] at (-2.05,0) {$270^\circ$};
  \node[anchor=north,align=center] at (0,-1.98)
    {$180^\circ$\quad\FigLang{직접 글린트 가지}{Direct-glint branch}};
  \fill[ocrtNavy] (0,0) circle (2pt);
  \node[note,anchor=north] at (0,-2.6)
    {\FigLang{거울 방위 $20^\circ\leftrightarrow340^\circ$: 같은 $I,Q$, 반대 $U$}{Mirror pair $20^\circ\leftrightarrow340^\circ$: same $I,Q$, opposite $U$}};
\end{scope}
\end{tikzpicture}
\caption{\FigLang{위에서 본 수평 look 방위의 차이로 RAA를 정한다. 영점과 시계 방향은 이 그림에서 선택한 규약이다. 제품의 SAA/VAA를 같은 영점·회전 방향으로 맞춘 후 변환한다. 축대칭 조건에서 $20^\circ$와 $340^\circ$를 모두 $20^\circ$로 접으면 signed $U$에 필요한 거울 방향 정보가 사라진다. 실제 편광 비교에는 Stokes 기준면도 맞춘다.}{RAA is the difference between horizontal look azimuths viewed from above. The zero reference and clockwise rotation are choices declared for this figure. Align the product's SAA/VAA origins and rotation directions before conversion. Under axisymmetric conditions, folding both $20^\circ$ and $340^\circ$ to $20^\circ$ discards the mirror-direction information needed for signed $U$. Polarization comparisons also require aligned Stokes reference planes.}}
\label{fig:geometry-azimuth}
\end{figure}


\subsection{When only RAA folded into 0--180\textdegree{} is available}
\label{sec:auto:01_concepts_geometry:4.2}
A folded value is commonly generated as follows.


\begin{ManualCode}
raa_360 = (vaa_deg - saa_deg) % 360.0
raa_folded = min(raa_360, 360.0 - raa_360)
\end{ManualCode}

Under axisymmetric conditions, a half-circle grid can be used for LUTs containing only \texttt{I} and \texttt{Q}. However, retaining only \texttt{raa\_\allowbreak{}folded=\allowbreak{}20\textdegree{}} loses whether the original direction was \texttt{20\textdegree{}} or \texttt{340\textdegree{}}. To preserve \texttt{U}, also store the original 0--360\textdegree{} direction or an indication of the mirror direction. Storing the absolute value \texttt{abs(U)} discards directional information.

\section{Matching Phi in other codes}
\label{sec:auto:01_concepts_geometry:5}
The current \texttt{rt\_\allowbreak{}raa\_\allowbreak{}convention.\allowbreak{}h} defines \textbf{one of the following two equivalent alternatives} for OCRT--OSOAA comparisons.


{\TableFont
\begin{longtable}{@{}>{\raggedright\arraybackslash}p{\dimexpr 0.3000\linewidth-5.33pt\relax}>{\raggedright\arraybackslash}p{\dimexpr 0.2300\linewidth-5.33pt\relax}>{\raggedright\arraybackslash}p{\dimexpr 0.4700\linewidth-5.33pt\relax}@{}}
\toprule
\textbf{Conversion method} & \textbf{Azimuth passed to OSOAA} & \textbf{U comparison with OCRT} \\
\midrule
\endfirsthead
\toprule
\textbf{Conversion method} & \textbf{Azimuth passed to OSOAA} & \textbf{U comparison with OCRT} \\
\midrule
\endhead
\bottomrule
\endfoot
A & \texttt{Phi\_\allowbreak{}OSOAA=\allowbreak{}(180-RAA\_\allowbreak{}OCRT) mod 360} & \texttt{U\_\allowbreak{}OCRT=\allowbreak{}-U\_\allowbreak{}OSOAA} \\
B & \texttt{Phi\_\allowbreak{}OSOAA=\allowbreak{}(RAA\_\allowbreak{}OCRT+180) mod 36\allowbreak{}0} & \texttt{U\_\allowbreak{}OCRT=\allowbreak{}U\_\allowbreak{}OSOAA} \\
\end{longtable}
}

Both methods construct the same physical comparison for \texttt{I/\allowbreak{}Q}. Do not apply method A's U sign reversal together with method B's angle transformation. This rule assumes the output conventions of the corresponding comparison harness; it is not a universal formula for arbitrary codes.

Current OCRT atmospheric outputs, above-surface \texttt{Rrs}, and below-surface \texttt{rrs} are Fourier-reconstructed using \textbf{the same public RAA}. The historical no-conversion rule \texttt{Phi\_\allowbreak{}OSOAA,\allowbreak{}water=\allowbreak{}RAA\_\allowbreak{}OCRT} corresponded to water outputs before the correction. Do not use it for water outputs from the current source.

When reading a CCRR reference CSV, \texttt{-\allowbreak{}-\allowbreak{}ccrr-\allowbreak{}compare} internally applies \texttt{RAA\_\allowbreak{}OCRT=\allowbreak{}(180-RAA\_\allowbreak{}CCRR) mod 360}. If the CSV's \texttt{raa\_\allowbreak{}deg} already follows the OCRT convention, do not feed it directly into this comparison mode and transform it twice. Check \texttt{raa\_\allowbreak{}ccrr\_\allowbreak{}deg}, \texttt{raa\_\allowbreak{}ocrt\_\allowbreak{}deg}, and \texttt{raa\_\allowbreak{}mapping} in the results.

\section{Stokes signs and off-principal-plane validation}
\label{sec:auto:01_concepts_geometry:6}
The public default path uses Stokes components aligned with the viewing-direction meridian plane (the plane containing the viewing direction and the vertical). The default surface Q convention is \texttt{q\_\allowbreak{}convention=\allowbreak{}1}, with the Fresnel core component \texttt{Q\_\allowbreak{}kernel=\allowbreak{}(Rp-Rs)/\allowbreak{}2}. \texttt{Q} and \texttt{U} are signed quantities, and negative values are normal. Comparison with sensor coordinates or scattering-plane coordinates may require a reference-plane rotation.

For an axisymmetric surface and medium, using the same reference plane, the following symmetries hold.


\begin{ManualCode}
I(RAA) =  I(360−RAA)
Q(RAA) =  Q(360−RAA)
U(RAA) = −U(360−RAA)
\end{ManualCode}

In the principal plane at \texttt{RAA=\allowbreak{}0/\allowbreak{}180\textdegree{}}, U is zero or at the numerical-error level, so errors in the U sign are not exposed. Also check off-principal-plane pairs such as \texttt{RAA=\allowbreak{}45/\allowbreak{}315\textdegree{}}. At exactly \texttt{VZA=\allowbreak{}0\textdegree{}}, the meridian plane is geometrically degenerate, so Q/U as a function of RAA may not correspond to intuitive left/right directions. Include small positive VZA values in model comparisons rather than using only nadir values.

\section{Scattering angle and conflicting historical comments}
\label{sec:auto:01_concepts_geometry:7}
Taking the dot product of the actual solar propagation vector and the upward viewing vector derived from the look directions declared above gives the atmospheric scattering angle:


\begin{ManualCode}
mu_s = cos(SZA), mu_v = cos(VZA)
cos(Theta) = −mu_s*mu_v − sin(SZA)*sin(VZA)*cos(RAA_OCRT)
\end{ManualCode}

This matches the expression currently used by \texttt{rt\_\allowbreak{}ipss\_\allowbreak{}cos\_\allowbreak{}scattering\_\allowbreak{}angle()}. For \texttt{SZA=\allowbreak{}VZA=\allowbreak{}30\textdegree{}}, \texttt{Theta=\allowbreak{}180\textdegree{}} at \texttt{RAA=\allowbreak{}0\textdegree{}}, whereas \texttt{Theta=\allowbreak{}120\textdegree{}} at the glint branch, \texttt{RAA=\allowbreak{}180\textdegree{}}. \textbf{The surface specular-reflection direction and exact backscattering by atmospheric particles are different concepts.}

The current repository also contains an unresolved inconsistency in its internal descriptions. The helper \texttt{rt\_\allowbreak{}atm\_\allowbreak{}scatter\_\allowbreak{}cos\_\allowbreak{}from\_\allowbreak{}public\_\allowbreak{}raa()} in \texttt{rt\_\allowbreak{}raa\_\allowbreak{}convention.\allowbreak{}h} uses \texttt{+sin(SZA)*sin(VZ\allowbreak{}A)*cos(RAA)}, opposite to the sign above, and associated comments and \texttt{test\_\allowbreak{}raa\_\allowbreak{}convention.\allowbreak{}c} describe \texttt{RAA=\allowbreak{}180\textdegree{}} as exact backscattering. At the time this manual was prepared, no production \texttt{src} calls to this helper were found; actual Fourier reconstruction uses \texttt{public\_\allowbreak{}RAA+\ensuremath{\pi}}, and IPSS uses the \texttt{-cos(RAA)} expression above. Therefore, do not use this helper or the “exact backscattering” wording in historical audit reports as the basis for satellite-geometry conversion. This manual does not change the source code.

\subsection{Implementation references}
\label{sec:auto:01_concepts_geometry:7.1}
\begin{itemize}
\item \href{https://github.com/brtnt/OCRT/blob/54861c68e35ec5aaa3938fad53e93dc2dfa7eb1d/code/OCRT_C/src/rt_raa_convention.h}{Public RAA helpers and current OSOAA mapping}
\item \href{https://github.com/brtnt/OCRT/blob/54861c68e35ec5aaa3938fad53e93dc2dfa7eb1d/code/OCRT_C/src/rt_fourier.c}{Internal +\ensuremath{\pi} conversion in Fourier reconstruction}
\item \href{https://github.com/brtnt/OCRT/blob/54861c68e35ec5aaa3938fad53e93dc2dfa7eb1d/code/OCRT_C/src/rt_ipss.c}{IPSS scattering angle and pixel-referenced geometry}
\item \href{https://github.com/brtnt/OCRT/blob/54861c68e35ec5aaa3938fad53e93dc2dfa7eb1d/code/OCRT_C/src/main.c}{CLI mode selection and CCRR conversion}
\item \href{https://github.com/brtnt/OCRT/blob/54861c68e35ec5aaa3938fad53e93dc2dfa7eb1d/code/OCRT_C/src/shared/surface.c}{Surface Stokes components and meridian-plane rotation}
\end{itemize}

Before changing input parameters, first align these three features: glint at \texttt{RAA=\allowbreak{}180\textdegree{}}, the opposite branch at \texttt{0\textdegree{}}, and the off-principal-plane U sign. This reduces common 180\textdegree{} and sign errors in comparisons.

